\documentclass[11pt]{article}

% =========================
% Packages (Minimal, Stable)
% =========================
\usepackage{amsmath, amssymb}
\usepackage{geometry}
\usepackage{hyperref}
\usepackage{enumitem}
\usepackage{longtable}

\geometry{margin=1in}

% =========================
% Metadata
% =========================
\title{Kernel L0 Dashboard \\ Mathematical and Accounting Specification}
\author{Aequitas}
\date{Version 1.0 --- Frozen Specification}

% =========================
% Custom Commands
% =========================
\newcommand{\AccountSet}[1]{\mathcal{A}_{#1}}
\newcommand{\Balance}{B}
\newcommand{\Total}{T}
\newcommand{\Period}{p}
\newcommand{\Prior}{p^{-}}
\newcommand{\Real}{\mathbb{R}}

% =========================
% Document
% =========================
\begin{document}
\maketitle

\section*{Normative Status}
This document is a \textbf{normative specification}.
It defines the authoritative mathematical and accounting meaning
of Kernel L0 Dashboard Tier~1 metrics.

Any implementation, test, or visualization that contradicts this
document is incorrect by definition.

\section*{Specification Mode}
Each metric is defined simultaneously in:
\begin{enumerate}
  \item Formal Mathematical Notation
  \item Accounting Narrative Language (GAAP-style)
\end{enumerate}

Both representations are \textbf{equally authoritative}.
If a discrepancy exists, the specification is invalid and must be
revised via a delta document.

\section*{Scope and Constraints}
\begin{itemize}
  \item Accrual basis
  \item Kernel L0 accounts only
  \item Read-only (no journal entry creation)
  \item No projections, recommendations, or scoring
  \item Period-aware
\end{itemize}

% =========================================================
\section{Preliminaries}
% =========================================================

\subsection{Fiscal Period Model}
Let $\Period$ denote a fiscal period selected by the user
(e.g., month, quarter, year).

Let $\Prior$ denote the immediately preceding fiscal period.

\subsection{Measurement Functions}

Because balance sheet accounts are \emph{stock} measures and
income statement accounts are \emph{flow} measures, define:

\begin{itemize}
  \item Period-end balance function:
  \[
    \Balance(a, \Period) \in \Real
  \]
  representing the balance (in reporting polarity) of account $a$
  at the close of period $\Period$.

  \item Period activity function:
  \[
    \Total(a, \Period) \in \Real
  \]
  representing the total activity (in reporting polarity) of
  income statement account $a$ over period $\Period$.
\end{itemize}

\subsection{Delta Operator}
For any metric $M(\Period)$:
\[
\Delta M(\Period) = M(\Period) - M(\Prior)
\]

If either term is undefined, $\Delta M(\Period)$ is undefined.

% =========================================================
\section{Kernel L0 Account Sets}
% =========================================================

\subsection{Cash}
\[
\AccountSet{cash} = \{10000, 10100\}
\]

\subsection{Revenue and Refunds}
\[
\AccountSet{rev} = \{40000\}
\qquad
\AccountSet{ref} = \{49000\}
\]

\subsection{Operating Costs}
\[
\AccountSet{cogs} = \{50000\}
\]
\[
\AccountSet{opex} = \{60000\}
\]
\[
\AccountSet{pay} = \{61000\}
\]
\[
\AccountSet{dep} = \{62000\}
\]

\subsection{Current Assets (Kernel Approximation)}
\[
\AccountSet{CA} = \{10000, 10100, 12000, 14000\}
\]

\subsection{Current Liabilities (Kernel Approximation)}
\[
\AccountSet{CL} = \{20000, 21000, 22000, 23000\}
\]

% =========================================================
\section{Metric Specifications}
% =========================================================

% =========================
\subsection{Metric: Total Cash}
% =========================

\subsubsection*{Formal Mathematical Definition}
\[
\text{TotalCash}(\Period)
= \sum_{a \in \AccountSet{cash}} \Balance(a, \Period)
\]

Domain:
\[
\text{TotalCash}(\Period) \in \Real
\]

\subsubsection*{Accounting Narrative Definition}
Total Cash represents immediate liquidity available to the entity,
comprising operating cash and undeposited funds.
It is derived from period-end balances.

\subsubsection*{Undefined Cases}
None.

\subsubsection*{Audit Rule}
Reproducible from the trial balance without estimates.

% =========================
\subsection{Metric: Net Revenue}
% =========================

\subsubsection*{Formal Mathematical Definition}
\[
\text{NetRevenue}(\Period)
= \sum_{a \in \AccountSet{rev}} \Total(a, \Period)
- \sum_{a \in \AccountSet{ref}} \Total(a, \Period)
\]

\subsubsection*{Accounting Narrative Definition}
Net Revenue represents recognized revenue reduced by refunds and
allowances for the selected period.

\subsubsection*{Undefined Cases}
None.

\subsubsection*{Audit Rule}
Derived solely from income statement totals.

% =========================
\subsection{Metric: Operating Income}
% =========================

\subsubsection*{Formal Mathematical Definition}

Let:
\[
\text{OperatingCosts}(\Period)
=
\sum_{a \in \AccountSet{cogs} \cup \AccountSet{opex} \cup
           \AccountSet{pay} \cup \AccountSet{dep}}
\Total(a, \Period)
\]

Then:
\[
\text{OperatingIncome}(\Period)
=
\text{NetRevenue}(\Period) - \text{OperatingCosts}(\Period)
\]

\subsubsection*{Accounting Narrative Definition}
Operating Income represents profit from core operations before
taxes, interest, and financing effects.

\subsubsection*{Explicit Exclusions}
Taxes, interest, equity effects, financing activities.

\subsubsection*{Undefined Cases}
None.

% =========================
\subsection{Metric: Net Working Capital}
% =========================

\subsubsection*{Formal Mathematical Definition}

\[
\text{CurrentAssets}(\Period)
= \sum_{a \in \AccountSet{CA}} \Balance(a, \Period)
\]

\[
\text{CurrentLiabilities}(\Period)
= \sum_{a \in \AccountSet{CL}} \Balance(a, \Period)
\]

\[
\text{NWC}(\Period)
=
\text{CurrentAssets}(\Period)
-
\text{CurrentLiabilities}(\Period)
\]

\subsubsection*{Accounting Narrative Definition}
Net Working Capital represents the excess of current assets over
current liabilities using the Kernel-defined approximation.

\subsubsection*{Undefined Cases}
None.

% =========================
\subsection{Metric: Current Ratio}
% =========================

\subsubsection*{Formal Mathematical Definition}
\[
\text{CurrentRatio}(\Period)
=
\begin{cases}
\dfrac{\text{CurrentAssets}(\Period)}
      {\text{CurrentLiabilities}(\Period)},
& \text{if } \text{CurrentLiabilities}(\Period) \neq 0
\\[1em]
\text{undefined},
& \text{if } \text{CurrentLiabilities}(\Period) = 0
\end{cases}
\]

\subsubsection*{Accounting Narrative Definition}
The Current Ratio is a liquidity indicator expressing current assets
relative to current liabilities.

\subsubsection*{Undefined Cases}
Division by zero; rendered as N/A.

% =========================================================
\section*{Freeze Declaration}
This specification is frozen for Kernel Dashboard v1.0.
Any modification requires a delta document and version increment.

% =========================================================
\appendix
\section{Worked Numerical Examples}
% =========================================================

This appendix provides concrete numerical instantiations of the
formal definitions in Section~3.
All examples are authoritative and must reconcile exactly in any
implementation.

All balances and totals are expressed in reporting polarity.

% =========================
\subsection{Example A.1 — Empty Company (Baseline)}
% =========================

\paragraph{Inputs (Period $\Period$)}
For all Kernel L0 accounts $a$:
\[
\Balance(a,\Period) = 0
\quad\text{and}\quad
\Total(a,\Period) = 0
\]

\paragraph{Results}
\[
\text{TotalCash}(\Period) = 0
\]
\[
\text{NetRevenue}(\Period) = 0
\]
\[
\text{OperatingIncome}(\Period) = 0
\]
\[
\text{NWC}(\Period) = 0
\]
\[
\text{CurrentRatio}(\Period) = \text{undefined}
\]

\paragraph{Interpretation Constraint}
Zero is a valid value.
Undefined ratios must be rendered as \emph{N/A}.

% =========================
\subsection{Example A.2 — Cash Only}
% =========================

\paragraph{Inputs}
\[
\Balance(10000,\Period) = 10{,}000
\]
All other balances and totals are zero.

\paragraph{Results}
\[
\text{TotalCash}(\Period) = 10{,}000
\]
\[
\text{NWC}(\Period) = 10{,}000
\]
\[
\text{CurrentRatio}(\Period) = \text{undefined}
\]

% =========================
\subsection{Example A.3 — Revenue with Refunds}
% =========================

\paragraph{Inputs}
\[
\Total(40000,\Period) = 50{,}000
\quad
\Total(49000,\Period) = 8{,}000
\]

\paragraph{Results}
\[
\text{NetRevenue}(\Period)
=
50{,}000 - 8{,}000
=
42{,}000
\]

\[
\text{OperatingIncome}(\Period) = 42{,}000
\]

% =========================
\subsection{Example A.4 — Full Operating Stack (Profit)}
% =========================

\paragraph{Inputs}
\[
\Total(40000,\Period) = 100{,}000
\]
\[
\Total(49000,\Period) = 5{,}000
\]
\[
\Total(50000,\Period) = 40{,}000
\]
\[
\Total(60000,\Period) = 20{,}000
\]
\[
\Total(61000,\Period) = 25{,}000
\]
\[
\Total(62000,\Period) = 3{,}000
\]

\paragraph{Intermediate}
\[
\text{NetRevenue}(\Period) = 95{,}000
\]
\[
\text{OperatingCosts}(\Period) = 88{,}000
\]

\paragraph{Results}
\[
\text{OperatingIncome}(\Period) = 7{,}000
\]

% =========================
\subsection{Example A.5 — Net Working Capital (Positive)}
% =========================

\paragraph{Inputs — Assets}
\[
\Balance(10000,\Period) = 20{,}000
\]
\[
\Balance(10100,\Period) = 500
\]
\[
\Balance(12000,\Period) = 12{,}000
\]
\[
\Balance(14000,\Period) = 1{,}500
\]

\paragraph{Inputs — Liabilities}
\[
\Balance(20000,\Period) = 8{,}000
\]
\[
\Balance(21000,\Period) = 2{,}000
\]
\[
\Balance(22000,\Period) = 3{,}000
\]
\[
\Balance(23000,\Period) = 6{,}000
\]

\paragraph{Intermediate}
\[
\text{CurrentAssets}(\Period) = 34{,}000
\]
\[
\text{CurrentLiabilities}(\Period) = 19{,}000
\]

\paragraph{Results}
\[
\text{NWC}(\Period) = 15{,}000
\]
\[
\text{CurrentRatio}(\Period) =
\frac{34{,}000}{19{,}000}
\approx 1.78947
\]

% =========================
\subsection{Example A.6 — Division by Zero (Ratio Undefined)}
% =========================

\paragraph{Inputs}
\[
\text{CurrentAssets}(\Period) = 5{,}000
\quad
\text{CurrentLiabilities}(\Period) = 0
\]

\paragraph{Results}
\[
\text{NWC}(\Period) = 5{,}000
\]
\[
\text{CurrentRatio}(\Period) = \text{undefined}
\]

\paragraph{Rendering Rule}
Undefined ratios must not be coerced into numeric values.

% =========================
\subsection{Example A.7 — Delta Computation}
% =========================

\paragraph{Inputs}
\[
\text{TotalCash}(\Prior) = 9{,}000
\quad
\text{TotalCash}(\Period) = 13{,}700
\]

\paragraph{Result}
\[
\Delta \text{TotalCash}(\Period)
=
13{,}700 - 9{,}000
=
4{,}700
\]

\paragraph{Delta Constraint}
If either operand is undefined, the delta is undefined.

% =========================================================
\section{Audit Reconciliation Procedure}
% =========================================================

For any metric defined in this specification, an independent reviewer
must be able to:

\begin{enumerate}[label=\arabic*.]
  \item Extract the relevant account balances and totals from the trial balance.
  \item Identify the applicable account sets defined in Section~2.
  \item Apply the formal mathematical definitions exactly as written.
  \item Reconcile the computed value to the dashboard output.
\end{enumerate}

Failure at any step indicates non-compliance with this specification.

\end{document}
